% !TeX root = ../main.tex
\chapter{绪论}

\section{研究背景与意义}

\subsection{XML 数据的特征与数据库应用背景}

可扩展标记语言（Extensible Markup Language，XML）通过元素及其嵌套关系表达层次化数据，长期用于跨系统数据交换、业务文档、配置数据和异构系统之间的信息传递。XML 1.0 规定了 XML 文档的标记语法与基本良构性要求，但不规定数据库必须采用何种物理存储或索引方式[1]。XML 的自描述特征使数据结构可以随应用变化，同时也使标签名称、层次路径和文本值在同一文档或多个文档之间反复出现。

与关系型表中固定模式的记录相比，XML 文档具有明显的层次性、顺序性和结构可变性。元素的嵌套关系、节点顺序、空元素以及文本内容共同构成文档语义。将 XML 放入关系型数据库统一管理时，数据库不仅需要保存文档内容，还需要在存储、查询、索引和类型转换之间维持一致语义。因此，数据库 XML 类型的物理表示不能只关注字节数，还需要考虑压缩之后的结构和值信息是否仍然能够参与查询。

PostgreSQL 为 XML 类型提供 \texttt{xpath}、\texttt{xpath\_exists} 等查询函数，并通过 libxml2 执行 XPath 1.0 表达式[2][3]。这为数据库 XML 查询提供了较为成熟的实现基线，也为研究压缩表示如何接入数据库类型和索引机制提供了可扩展的平台。本文以 PostgreSQL 的 XPath 1.0 查询语境为应用背景，围绕路径存在、路径值获取和精确内容匹配等查询意图，研究 XML 压缩表示与查询、索引处理之间的协同关系。

\subsection{XML 存储与查询的现实需求}

XML 存储与查询之间存在若干相互关联的需求。首先，XML 的层次结构不能简单等同于无序的键值集合。将 XML 压平为若干路径和值，便于去重和查询，但如果没有额外结构信息，可能无法准确恢复父子边界、重复容器、节点顺序或空元素。其次，压缩和查询可能重复处理同一语义信息。写入时解析器已经识别了元素路径和文本值；如果查询时仍从原始 XML 重新解析这些信息，压缩阶段形成的中间结果就没有被查询复用；如果在原始 XML 之外另行物化查询数据，又会增加派生数据、索引维护和写入成本。

不同查询所需的信息也并不相同。在 XPath 1.0 语境下，路径存在性需要判断某条路径是否出现；路径值获取需要沿文档顺序找到匹配路径对应的值；内容查询则需要对文本值进行精确判断。因而，XML 压缩方法需要明确编码表示能够直接提供哪些信息、查询过程需要解码哪些信息，以及候选结果如何进行精确验证。索引优化还要同时考虑候选生成和结果正确性：Bloom Filter 可以快速排除部分不可能的候选，但不能把“可能存在”直接当作真实命中；哈希标识也可能发生碰撞，必要时必须回到压缩值或原始值进行精确复核。

\subsection{研究意义}

本文研究如何复用 XML 压缩编码过程中形成的路径标识、值标识和结构记录，使它们在压缩表示生成后继续服务于路径判断、路径值获取、内容匹配及索引键提取。本文强调的不是将 XML 当作普通文本进行压缩，而是利用 XML 的结构特征识别重复路径、重复值和节点关系，形成结构感知的压缩表示；在此基础上，使编码阶段已经产生的中间数据继续被后续查询和索引处理使用。

本研究的意义主要体现在两个方面。方法层面，建立从 XML 解析、结构感知编码、中间数据组织到查询与索引复用的连续处理过程，分析同一组编码信息如何减少后续查询中的重复识别，并明确候选过滤、精确验证和结构恢复之间的正确性关系。应用层面，在 PostgreSQL 扩展原型中将该方法落实为二进制表示、XPath 导向查询函数和 GIN 索引处理流程，再通过实验观察其存储开销、查询表现及写入代价。PostgreSQL 为方法提供 XML 查询基线和可扩展索引环境，使本文能够在数据库类型层面对方法的可行性和基本能力进行验证。

\section{国内外研究现状}

\subsection{XML 压缩方法研究现状}

XML 文档包含重复出现的标签、路径及文本值，相关压缩研究因此从通用文本压缩逐步发展出利用 XML 结构特征的表示方法。XMill 将 XML 内容按结构特征分组，再结合通用压缩技术处理同类数据，体现了结构信息参与压缩组织的思路[5]。国内 XSC 等工作从 XML 数据流和结构事件出发组织编码，并关注压缩数据的访问与 XPath 处理[4]。这类方法表明，XML 压缩并非只能把文档视为普通字符序列；元素路径、节点类别和文本内容都可以影响数据组织及编码方式。

从已有研究可以看到，XML 压缩通常需要在压缩率、解码粒度和结构保真之间进行权衡。面向压缩率的表示可以通过字典、统计编码或通用压缩器减少字节数，但查询时可能需要恢复较大范围的数据；结构感知方法能够利用路径和节点关系减少结构性冗余，但需要额外保存解码和导航所需的结构信息。字典作用域、编码单位、结构事件表示和解码粒度都会影响最终的空间与访问成本。

\subsection{可查询 XML 压缩研究现状}

随着数据库对压缩数据直接访问的需求增加，研究进一步形成了可查询 XML 压缩路线。XGrind 和 XPRESS 研究在压缩表示上进行内容或结构查询的方法[6][7]；Buneman、Grohe 与 Koch 讨论在压缩 XML 骨架上执行路径查询的方法[8]。这些工作表明，压缩表示可以保留或构造查询可使用的结构与内容信息，查询不必在所有情况下都先完整恢复原始文档。

后续研究进一步把查询负载和访问粒度纳入压缩数据组织。XQueC 根据查询需求选择结构与数据容器的组织方式[9]；XSAQCT 关注 XML 压缩数据的随机访问及查询能力[10]；XQzip 利用结构索引辅助压缩 XML 查询[11]；XBzipIndex 研究压缩 XML 的导航和搜索[12]；XXS 与 SXSI 则探索压缩文档或自索引结构上的 XPath 查询[13][14]。这些方法在支持的查询语言、编码格式、索引形式、压缩单位及查询时解码范围方面各有特点，共同推动了 XML 压缩研究从“压得更小”扩展到“压缩后如何定位并处理查询所需的信息”。

\subsection{XML 查询与索引研究现状}

XML 查询研究包括对文档树进行路径导航、对节点内容进行匹配，以及借助索引减少候选范围等方向。DataGuides 使用结构摘要描述半结构化数据中的实际路径[15]；IndexFabric 将路径组织为索引键以支持半结构化数据检索[16]；XISS/R 则研究在关系数据库环境中存储与索引 XML[17]。针对压缩 XML，XQzip、XBzipIndex、XXS 和 SXSI 等工作把结构访问与压缩表示或索引组织结合起来[11]--[14]。

在数据库实现方面，PostgreSQL 提供 XML 类型以及基于 libxml2 的 XPath 1.0 查询函数[2][3]，并通过扩展索引接口支持不同数据类型定义键抽取、候选匹配和结果复核机制[19]。GIN（Generalized Inverted Index，通用倒排索引）可以将一个复合值拆解为多个索引键，在查询阶段先定位候选行，再由一致性函数和执行器完成必要的复核。已有研究和系统共同说明，XML 查询可以由文档内结构访问、跨文档索引和数据库执行器协作完成；同时，索引还会带来额外的存储、构建和维护成本。

\subsection{研究现状小结}

总体来看，国内外研究已经形成结构感知压缩、可查询压缩、路径索引、内容检索和数据库集成等多条技术路线。现有工作为本文提供了重要基础：XML 的结构和内容可以在编码阶段被分别识别和组织，压缩表示可以保留查询所需的信息，结构索引和倒排索引可以减少无关访问，数据库扩展接口可以将自定义表示接入查询执行过程。

在这些研究基础上，本文进一步关注编码阶段信息在多个后续环节之间的复用关系：路径和值在编码时被识别之后，能否以一致表示同时服务于压缩存储、XPath 导向查询和索引键提取；为此保留的结构和辅助信息会带来怎样的空间及处理代价；Bloom Filter 或哈希候选如何与精确验证结合以保持查询结果正确。上述问题构成本文研究目标与研究内容的出发点。

\section{研究目标与研究内容}

\subsection{研究目标}

本文的总体目标是面向数据库 XML 类型，研究一种以结构感知压缩和压缩编码复用为核心的压缩与索引优化方法。该方法在 XML 解析阶段识别元素路径、文本值和节点结构，将可重复的路径和值映射为可复用标识，并组织为可解码的二进制压缩表示；在查询和索引处理阶段继续消费这些编码标识，从而减少对同一 XML 语义信息的重复解析和重复识别。

围绕总体目标，本文重点解决以下研究问题：

\begin{enumerate}[label=（\arabic*）]
\item 如何构造同时保留 XML 结构、文档顺序和内容信息的压缩表示，使结构感知压缩不会破坏节点语义和解码正确性；
\item 如何使压缩编码阶段形成的路径 ID、文本 token ID 和结构记录被压缩表示、XPath 导向查询函数及索引键提取共同复用；
\item 如何将路径存在、路径值获取和精确内容匹配等查询过程分解为候选定位与精确验证，并明确 Bloom Filter、路径索引和 GIN 索引各自的作用范围；
\item 在统一统计口径下，如何分别评价压缩列值、完整 relation、索引和共享字典的空间开销，以及查询性能和写入代价。
\end{enumerate}

本文以 PostgreSQL 的 XML 类型和 XPath 1.0 查询能力作为应用语境和实现基线。原型围绕上述查询意图实现相应函数和索引接口，研究结论限定在当前方法、原型、数据集和测试配置范围内。

\subsection{研究内容}

围绕上述目标，本文开展以下四项相互衔接的研究工作。

\begin{enumerate}[label=（\arabic*）]
\item \textbf{压缩表示与查询操作的语义定义。} 对 XML 解析所得的路径、叶文本、节点顺序和结构关系进行形式化表示，明确路径存在、路径值获取和精确内容匹配等查询操作的输入、输出及正确性条件。
\item \textbf{面向查询的结构感知压缩方法。} 复用编码阶段形成的路径 ID 和文本 token ID，将其组织为文档序内容流、排序路径索引、Bloom 元数据、长值内联区和结构流；短值使用字典标识，整数使用 Varint 编码，同时保留恢复 XML 结构所需的信息。
\item \textbf{基于压缩表示的查询与索引方法。} 在文档内部利用路径索引和内容流完成路径判断、路径值获取及精确内容查询，在关系级利用 GIN 生成候选行；对于 Bloom Filter 假阳性和长值哈希碰撞，均通过精确扫描或谓词 recheck 保证结果语义。
\item \textbf{原型实现与实验评估。} 在 PostgreSQL 扩展中实现自定义数据类型、输入输出函数、查询函数和 GIN 操作类，通过编码/解码正确性测试、空间测量、查询性能测试和消融观察评估方法的可行性与代价边界。
\end{enumerate}

\subsection{主要贡献}

本文的主要贡献可以概括为以下三点：

\begin{enumerate}[label=（\arabic*）]
\item 提出以压缩编码中间数据复用为核心的 XML 压缩与索引优化方法，将编码阶段形成的路径和值标识组织为可被压缩存储、查询处理和索引键提取共同消费的中间表示；
\item 建立结构感知压缩表示与查询处理之间的对应关系，在保留文档序内容和节点结构的同时，结合路径索引、Bloom Filter、GIN 候选定位及精确验证支持 XPath 导向查询；
\item 在 PostgreSQL 扩展原型中验证方法的可实现性，并从列值、relation、索引、共享字典、查询和写入等多个维度报告实验观察，明确不同统计口径和实验条件下的适用边界。
\end{enumerate}

本文的创新主张针对上述方法流程及其正确性约束，不将字典编码、Varint、Bloom Filter 或倒排索引单独宣称为新算法，也不把当前原型实验结果泛化为对其他数据库平台或所有可查询 XML 压缩系统的普遍性能结论。

\subsection{研究范围与设计要求}

本文研究范围是数据库 XML 类型的压缩表示、查询处理和索引优化方法，重点考察结构感知压缩和编码中间数据复用。方法设计遵循以下要求：一是利用 XML 的结构和内容特征进行压缩，避免把 XML 仅作为普通文本字节串处理；二是使编码阶段形成的路径和值标识供压缩、查询和索引键提取复用；三是保留文档顺序和结构信息，保证解码和查询语义；四是区分候选过滤与精确结果判定，避免 Bloom Filter 假阳性和哈希碰撞造成错误命中；五是分别报告空间、查询和写入代价，不用单一指标替代完整评价。

\section{论文组织结构}

\subsection{技术路线}

本文的技术路线如下：首先，从 XML 数据结构和数据库查询需求出发，分析结构性冗余、压缩后查询和索引维护之间的关系，确定编码中间数据复用这一核心问题；其次，对 XML 文档进行语义解析，提取路径、文本值、节点顺序和结构关系，并通过字典及变长整数编码形成结构感知的压缩表示；再次，利用同一组路径和值标识构造路径索引、Bloom 摘要和 GIN 索引键，使查询处理能够直接消费编码阶段形成的信息；最后，通过正确性测试以及存储、查询、写入和消融实验评价方法能力，分析方法的适用条件与有效性威胁。

\subsection{章节安排}

全文共分为七章，各章内容如下：

\begin{enumerate}[label=第\arabic*章]
\item 绪论。介绍 XML 数据库应用背景、XML 存储与查询需求、国内外研究现状，提出本文研究目标、研究内容、主要贡献和技术路线。
\item XML 压缩存储与查询的理论基础。介绍 XML 数据模型、结构感知压缩、可查询 XML 压缩、路径与内容索引、字典编码、Varint、Bloom Filter 以及 PostgreSQL 扩展和 GIN 机制。
\item 可查询 XML 压缩方法及问题分析。基于前人研究分析编码阶段信息复用、查询候选生成与精确验证，以及压缩、查询和索引代价之间的耦合关系，明确本文的方法要求。
\item 面向查询的 XML 压缩表示方法。给出方法总体思路、XML 结构与内容编码、字典复用、变长整数编码、辅助数据组织，以及压缩表示的解码和正确性约束。
\item 基于压缩表示的查询处理方法。说明查询模型、精确内容查询、路径存在与值提取、路径索引与 Bloom 过滤、候选验证以及 PostgreSQL 操作符和 GIN 索引实现。
\item 实验评估与结果讨论。围绕研究问题设计实验，评估编码正确性、存储开销、查询性能、索引效果和写入代价，并讨论实验结果的适用范围与有效性威胁。
\item 总结与展望。总结本文工作和主要结论，说明研究不足，并提出后续研究方向。
\end{enumerate}

\subsection{本章小结}

本章从 XML 数据特征和数据库应用需求出发，说明了 XML 存储、压缩、查询和索引之间的关联，概述了国内外 XML 压缩、可查询压缩以及 XML 查询索引研究，并在此基础上提出本文的研究目标、研究内容和主要贡献。本文的核心研究主张是：以结构感知方式压缩 XML，并复用压缩编码阶段产生的路径、值和结构信息，使其继续服务于 XPath 导向查询和索引键提取。下一章介绍理解该方法所需的理论基础与代表性技术路线。
